% Abstract. Numbers are generated macros (report/generated/numbers.tex).
An idle superconducting qubit waiting during a neighbour's mid-circuit measurement loses
coherence to stochastic dephasing and to deterministic frequency shifts. We ask whether
model-based optimization of dynamical-decoupling (DD) pulse timings and phases with CMA-ES
outperforms strong canonical sequences. To answer it fairly, we pre-registered a synthetic
transmon-like model, seeds, baselines and a training-only baseline-selection rule, and scored
every arm on independent held-out noise. A finite-pulse Monte-Carlo simulator agrees with
filter-function theory in five predetermined configurations (largest $|z|=\GthreeMaxZ$).
Synthetic Ramsey/Hahn/CPMG spectroscopy recovers a predictive noise model (held-out RMSE
$\SpecGaussFullRmseRatio\times$ the shot-noise floor), but the spectrum is not uniquely
identified. In the pre-registered Scenario~A, no CMA-ES mode beat the selected XY4/XY8 baseline,
whose infidelity ($\BenchASixteenSelEps$ at $n=16$) is already at the $T_1$-plus-white-noise
floor ($\FloorA$). With a deterministic Stark-like shift and virtual-$Z$ correction, the only
consistent gain was small (ratio $\BenchBFourBothRatioLo$--$\BenchBFourBothRatioHi$ at $n=4$).
Optimizing on fitted models, or with more headroom, did not help, and training gains did not
generalize. All results are synthetic; no quantum hardware was used.
